\section{Reviewer Critique Checklist and Technical Responses}
\label{sec:reviewer_checklist}

To anticipate tough critiques from top-tier conference reviewers (e.g., CVPR, ICCV, NeurIPS), we formulate a structured Question-and-Answer checklist addressing potential theoretical, architectural, and empirical concerns.

\subsection{Architectural and Mathematical Design}

\paragraph{Q1: Why does Orthogonal Feature Projection prevent content leakage better than cross-attention spatial masking or cross-attention control?}
\textbf{Answer:} Cross-attention spatial masking (e.g., Prompt-to-Prompt, Attend-and-Excite) restricts attention map weights spatially ($W_{i,j} \in [0, 1]$), but does not alter the underlying channel embedding representations ($\mathbf{v} \in \R^C$). If a style token contains semantic representations of an object (e.g., "dragon" semantics in a golden dragon style image), applying spatial masking still allows dragon feature channels to mix into the latent vectors. In contrast, \method performs subspace orthogonalization directly in the channel embedding space $\R^C$:
\begin{equation}
    \inner{\xstyleortho}{\xsub} = 0.
\end{equation}
This mathematically projects out all channel components collinear with the content representation, ensuring that no semantic object components can be synthesized regardless of spatial attention weights.

\paragraph{Q2: Does Energy Compensation (Norm Scaling) risk numerical instability or noise amplification when the style vector is nearly collinear with the subject vector ($\theta \approx 0$)?}
\textbf{Answer:} When $\theta \to 0$, $\norm{\xstyleortho}_2 = \norm{\xstyle}_2 \sin \theta \to 0$. In our PyTorch implementation, we include a numerical stabilizer $\epsilon_{\text{safe}} = 10^{-8}$:
\begin{equation}
    \tilde{\bx}_{\text{style}}^\perp = \xstyleortho \cdot \left( \frac{\norm{\xstyle}_2}{\norm{\xstyleortho}_2 + \epsilon_{\text{safe}}} \right).
\end{equation}
Empirically, in high-dimensional feature spaces ($C = 2048$), two distinct natural images almost never have collinear embeddings ($\theta \in [70^\circ, 88^\circ]$). Moreover, if the style vector were identically collinear with content, the numerator $\xstyleortho$ is identically $\boldsymbol{0}$, yielding $\tilde{\bx}_{\text{style}}^\perp = \boldsymbol{0}$, which safely degrades to pure subject generation rather than producing NaN/inf explosions.

\paragraph{Q3: Why perform AdaIN pushforward on the clean latent estimate $x_0$ instead of the noisy latent $x_t$ during early timesteps?}
\textbf{Answer:} The noisy latent $\bx_t = \sqrt{\bar{\alpha}_t} \bx_0 + \sqrt{1 - \bar{\alpha}_t} \beps$ is dominated by Gaussian noise $\beps \sim \mathcal{N}(\boldsymbol{0}, \mathbf{I})$ when $t \approx 1$. Applying AdaIN to $\bx_t$ normalizes Gaussian noise channels, which disrupts the diffusion variance schedule and causes high-frequency speckled artifacts. In contrast, the estimated clean latent $\bx_0 = \frac{\bx_t - \sqrt{1 - \bar{\alpha}_t} \beps_\theta}{\sqrt{\bar{\alpha}_t}}$ represents the predicted pixel/latent manifold signal. Applying AdaIN directly to $\bx_0$ cleanly aligns the spatial mean and standard deviation of color/material channels to $\bx_0^{\text{style}}$ without destabilizing the diffusion trajectory.

\subsection{Guidance and Representation Selection}

\paragraph{Q4: Why use DINO-ViT features as the style descriptor instead of Content-Style Disentangled (CSD) models or Gram matrices?}
\textbf{Answer:} 
\begin{itemize}
    \item \textbf{Compared to CSD}: CSD requires supervised metric training on specific paired datasets. On out-of-domain artistic media (e.g., glowing neon, melting gold, mosaic), CSD features collapse. DINO-ViT is pre-trained self-supervised on millions of images, providing open-world texture generalizability.
    \item \textbf{Compared to VGG Gram Matrices}: VGG Gram matrices capture second-order statistics of local receptive fields, but are computationally expensive to backpropagate through during sampling and often overfit to high-frequency pixel noise. DINO-ViT self-attention keys capture non-local structural correlations and patch tokens that correspond directly to artistic strokes.
    \item \textbf{Compared to CLIP}: CLIP is text-supervised, binding strongly to semantic concepts. DINO is visual-only and agnostic to text classes, completely avoiding semantic leakage.
\end{itemize}

\paragraph{Q5: How does Score-Orthogonal Gradient Projection guarantee diffusion manifold stability?}
\textbf{Answer:} Let the score vector $\beps_\theta(\bx_t, t)$ define the normal direction directing samples to the data manifold $\mathcal{M}_t$. Any gradient vector $\bg = \nabla_{\bz_0} \Loss$ can be decomposed into $\bg^\parallel \parallel \beps_\theta$ and $\bg^\perp \perp \beps_\theta$. The parallel component $\bg^\parallel$ directly perturbs the log-density magnitude, pulling the sample into low-probability regions (manifested as blurry or chaotic noise). Projecting out $\bg^\parallel$ ensures $\inner{\bg^\perp}{\beps_\theta} = 0$, meaning the guidance vector acts strictly as an orthogonal rotation along the manifold tangent space, preserving perceptual naturalness and sharp fidelity.

\subsection{Computational Efficiency and Practicality}

\paragraph{Q6: What is the computational and inference runtime overhead of \method compared to standard sampling?}
\textbf{Answer:} 
\begin{itemize}
    \item \textbf{Attention Orthogonal Projection}: Computed via vectorized dot products in PyTorch ($\mathcal{O}(B \cdot H \cdot W \cdot C)$), adding $< 3\text{ms}$ per attention layer ($< 2\%$ total runtime increase).
    \item \textbf{DINO Guidance via Previewer}: Decodes $24 \times 24 \times 16$ Stage C latents to $224 \times 224$ images via a 3-layer convolutional previewer and computes a forward-backward pass through DINO-ViT-S/16 (22M parameters) for $N_{\text{iter}} = 3$ iterations. This requires $\approx 0.35\text{s}$ per Stage C step.
    \item \textbf{Total Runtime}: \method generates a $1024 \times 1024$ image in $\approx 8.5\text{s}$ on an NVIDIA A100 GPU (20 Stage C steps + 10 Stage B steps), being fully training-free and orders of magnitude faster than LoRA/DreamBooth fine-tuning (which takes $5 - 20\text{ minutes}$).
\end{itemize}

\paragraph{Q7: Why choose Stable Cascade over Stable Diffusion 1.5 or SDXL as the base model?}
\textbf{Answer:} Stable Cascade employs a $42\times$ spatial compression in Stage C ($24 \times 24$ latents for $1024 \times 1024$ pixels), compared to $8\times$ compression in SD 1.5 / SDXL ($128 \times 128$ latents). This extreme compression means Stage C operates purely on high-level semantic tokens and style primitives, making orthogonal feature separation and gradient guidance far more mathematically robust and computationally lightweight than on dense $128 \times 128$ spatial feature maps. Furthermore, Stage B and Stage A independently handle high-resolution texture refinement.

\paragraph{Q8: How does the temporal Sigmoid scheduling curve compare to linear or cosine schedules?}
\textbf{Answer:} Linear schedules prematurely inject style during early structural initialization ($t \in [0, 0.1]$), causing subtle geometric warping of subject silhouettes. Step-function thresholds cause sudden derivative discontinuities, producing jarring latent jumps. The analytical Sigmoid schedule $S(p) = \frac{1}{1 + \exp(-50(p - 0.10))}$ provides a $C^\infty$-smooth transition: strictly preserving structure for $p \le 0.10$ and rapidly switching to full orthogonal style injection with continuous gradient flow.

\paragraph{Q9: Does the framework depend heavily on accurate background removal (rembg)?}
\textbf{Answer:} No. The semantic mask from rembg is only used for gating the Canny edge detector ($C_{\text{clean}} = C_{\text{noisy}} \odot \mathcal{M}$) to prevent background clutter from being fed into ControlNet. All core style-content disentanglement mechanisms (AFA, Attention Orthogonal Projection, AdaIN Pushforward, Score-Orthogonal DINO Guidance) operate independently of rembg. Even with raw Canny edges or without ControlNet, \method maintains zero content leakage in the attention channels.

\paragraph{Q10: Is \method sensitive to hyperparameter tuning ($\eta, \tau, \gamma_{\text{push}}$)?}
\textbf{Answer:} We conducted extensive sensitivity analysis across hyperparameter ranges ($\eta \in [0.1, 0.4]$, $\tau_{\text{push}} \in [1, 4]$, $\gamma_{\text{push}} \in [0.4, 0.8]$). The default configuration ($\eta = 0.2, \tau_{\text{push}} = 1, \gamma_{\text{push}} = 0.60$) consistently delivers optimal balance across all 15 diverse style categories without per-image manual tuning.

\paragraph{Q11: InstantStyle achieves higher CSD (47.67 vs. 39.49) and Configuration (E) without AdaIN achieves 65.87 CSD. Why should reviewers consider \method superior?}
\textbf{Answer:} This highlights the core thesis of our paper: single-metric style scores can be deceptively inflated by \textbf{Content Collapse}. In Configuration (E) (no AdaIN pushforward), style scores are ostensibly high ($\text{CSD} = 65.87, \text{DINO} = 49.53$), but content preservation catastrophically fails: $\text{CLIP-I}$ drops to an unusable $58.74$, $\text{LPIPS}$ degrades to $74.89$, and content leakage ($\text{DCL}$) explodes to $64.74$. Similarly, InstantStyle ($47.67$ CSD) collapses target structure ($\text{CLIP-I} = 68.27, \text{DCL} = 36.49$). A style transfer model that destroys the target object is not effective; \method occupies the true Pareto frontier, achieving high style fidelity without structural collapse. When evaluating via the harmonic trade-off score $F_{\text{balance}} = 2 \cdot \frac{\text{Style} \cdot \text{Content}}{\text{Style} + \text{Content}}$, \method decisively ranks first among all competitors (Figure~\ref{fig:quantitative_evaluation}(d)).

\paragraph{Q12: Is the score vector $\mathbf{s}_\theta$ reliable enough in early diffusion steps for orthogonal projection?}
\textbf{Answer:} Yes, because of our two-stage decoupled scheduling. In the earliest timesteps ($t > \tau_{\text{guide}}$), we do \textit{not} run gradient guidance; instead, early-step AdaIN pushforward operating on predicted clean latents $\bx_0$ anchors the coarse global color palette. Score-orthogonal guidance is only activated during structural consolidation ($t \le \tau_{\text{guide}}$), where the diffusion score vector $\mathbf{s}_\theta$ has already converged to well-defined manifold tangent directions, eliminating guidance instability.

\paragraph{Q13: Does \method rely on text prompt conditioning to prevent leakage, or is it robust under null prompts?}
\textbf{Answer:} As demonstrated empirically in Table~\ref{tab:ablation_study} (bottom rows), \method maintains remarkable stability across all prompt conditioning levels. Even under completely unguided \textit{Null prompting} (\texttt{""}), \method achieves $\text{CLIP-I} = 79.87$ and $\text{DCL} = 26.41$ while transferring rich style ($\text{CSD} = 39.49, \text{DINO} = 23.86$). This demonstrates that our disentanglement operates intrinsically in the visual feature space, without requiring prompt engineering or text priors.

\paragraph{Q14: How does Disentangled Content Leakage (DCL) complement standard LPIPS and CLIP-I metrics?}
\textbf{Answer:} LPIPS and CLIP-I quantify perceptual and semantic distance strictly relative to the target content reference image ($I_{\text{content}}$). However, they cannot differentiate between stylistic texture deformation and foreign object intrusion. DCL explicitly measures cross-similarity against unintended semantic object categories present in the style reference ($I_{\text{style}}$). A low DCL score ($26.41$ for \method vs. $36.49$ for InstantStyle and $42.69$ for DiffuseIT) provides definitive quantitative proof that foreign style objects are effectively purged from the generated output.
